Dvě potenciálová minima představují možné polohy iontu, oddělené energetickou bariérou $W$.
\begin{itemize}
\item \textbf{Bez pole:} obě polohy jsou rovnocenné. Pravděpodobnosti přeskoku v obou směrech jsou stejné a jsou úměrné $e^{-W/(k_BT)}$; nevzniká usměrněný pohyb.
\item \textbf{S polem:} pro kladný iont se bariéra ve směru pole sníží na $W-\Delta W$, proti poli zvýší na $W+\Delta W$. Přeskok ve směru pole je pravděpodobnější, vzniká usměrněný pohyb iontů.
\end{itemize}
$k_B$ je Boltzmannova konstanta a $T$ absolutní teplota. Model slouží k vysvětlení \textbf{tepelně aktivované iontové vodivosti a iontové relaxační polarizace}.
